%!TEX encoding = UTF8
%!TEX root = 0-notes.tex

\chapter{Polynômes à coefficients réels}
\label{chap:polynomes}

\section{Introduction}

\dfn{polynôme à coefficients réels, fonction polynomiale}{
	Soit $d\in\N$. Un \emphindex{polynôme} à coefficients réels est une expression du type
		\[ a_0 + a_1x + a_2x^2 + \dots + a_d x^d, \]
	où $a_d \neq 0$, $d$ est le \emphindex{degré} du polynôme, les $a_i\in\K$ sont ses \emphindex{coefficients}, et $x$ est une abstraite (pas nécessairement un nombre réel !).
	
	Une fonction $f(x) = a_0 + a_1x + a_2x^2 + \dots + a_d x^d$ de $\R$ dans $\R$ est une \emphindex{fonction polynomiale}.
}{dfn:polynome}

\nt{
	Un polynôme n'a pas nécessairement vocation à être vu comme une fonction.
	L'objet $x$ doit être compris comme indépendant des nombres, mais aussi de ses puissances. 
	La définition précise d'indépendance fait l'objet de la définition \ref{dfn:indep-lin}.
	
	Nous verrons également qu'une fonction polynomiale peut prendre des valeurs non nécessairement réelles ou complexes, comme par exemple des matrices (chapitre \ref{chap:calcul-matriciel}).
}


\notation{
	L'ensemble des polynômes à coefficients réels est noté $\R[x]$.
	
	Pour $n\in\N$, l'ensemble des polynômes à coefficients dans $\R$ de degré inférieur ou égal à $n$ est noté $\R_n[x]$.
}

\nomen{
	Le polynôme $a_0 + a_1x + a_2x^2 + \dots + a_d x^d$ est \emphindex{polynôme unitaire} dès que son \emphindex{coefficient dominant} $a_d$ est égal à 1.
}

\notation{
	Le polynôme $a_0 + a_1x + a_2x^2 + \dots + a_d x^d$ est noté $a_i x^i$ en utilisant la \emphindex{convention de sommation d'Einstein} dans laquelle chaque expression dans laquelle un indice $i$ apparaît deux fois sous-entend sa sommation sur toutes les valeurs que prend $i$.
	
	Elle nous permettra de nommer les coefficients sans écrire le polynôme ni définir son degré.
}

\dfn{addition et multiplication des polynômes}{
	Pour $f(x) = f_i x^i$ et $g(x) = g_i x^i$,
	\begin{align*}
		(f+g)(x) = (f_i + g_i)x^i && \et && (f\cdot g)(x) = \left(\sum_{k=0}^{i} f_k g_{i-k}\right)x^i.
	\end{align*}
}{prop:Rx-ring}

\section{Indépendance linéaire}

\nomen{
	Soient $p_1, p_2, \dots, p_k$ des polynômes à coefficients réels ($k\in\N^*$).
	Une \emphindex{combinaison linéaire} est une expression de la forme
		\[ \alpha_1 p_1 + \alpha_2 p_2 + \dots + \alpha_k p_k, \]
	où les $\alpha_i$ sont des nombres réels pour $i=1, 2, \dots, k$.
	
	De plus, 
	\begin{enumerate}
		\item si $\alpha_1 p_1 + \alpha_2 p_2 + \dots + \alpha_k p_k = 0$, alors la combinaison est une \emphindex{combinaison linéaire nulle} ; et
		\item si chaque coefficient $\alpha_i$ est nul, la combinaison est une \emphindex{combinaison linéaire triviale}.
	\end{enumerate}
}

\exe{}{
	Montrer que le polynôme $x^2 + x + 1$ est combinaison linéaire des polynômes $x^2 + 1$ et $x+1$.
}{exe:comb-lin-poly}{

}

\exe{}{
	Montrer que le polynôme $x^2 - 5x + 6$ n'est pas combinaison linéaire des polynômes $x - 2$ et $x-3$.
}{exe:comb-lin-poly}{

}

\dfn{indépendance linéaire}{
	Soit $d\in\N$. 
	La famille $\{ 1, x, x^2, \dots, x^d \}$ est \emphindex{linéairement indépendante sur $\R$} ou \emphindex{libre} car dès que
		\[ \alpha_0 + \alpha_1 x + \alpha_2 x^2 + \dots + \alpha_d x^d = 0 \]
	pour certains nombres réels $\alpha_0,\alpha_1, \alpha_2 \dots, \alpha_d\in\R$, alors
		\[ \alpha_0 = \alpha_1 = \alpha_2 = \dots = \alpha_d = 0. \]
	Autrement dit, la seule \emphindex{combinaison linéaire nulle} est la \emphindex{combinaison linéaire triviale}.
}{dfn:indep-lin}

\nomen{
	Une famille de polynôme non linéairement indépendante est dite \emphindex{famille liée}.
}

\nt{
	Les polynômes $\{ 1 ; x ; x^2 ; \dots \}$ sont indépendants sur $\R$ par définition.
	Les fonctions polynômiales associées, elles, sont indépendantes par démonstration. 
	L'exercice \ref{exe:1x-indep} traite du cas $\{ 1 ; x \}$ et peut se généraliser.
	
	La distinction est subtile mais importante car, sur un ensemble fini, un polynôme non nul peut être associé à une fonction nulle.
	L'exercice \ref{exe:fonction-nulle-S} fournit un exemple.
}

\nomen{
	La \emphindex{fonction nulle} est la fonction $f$ vérifiant $f(x) = 0$ pour tout $x\in\R$.
}

\exe{difficulty=1}{
	Montrer que la famille de fonctions polynomiales sur $\R$ $\{1 ; x\}$ est linéairement indépendante sur $\R$. 
	C'est-à-dire que dès que $\alpha_0 + \alpha_1 x$ est la fonction nulle sur $\R$ pour certains nombres réels $\alpha_0, \alpha_1 \in\R$, alors $\alpha_0 = \alpha_1 = 0$.
}{exe:1x-indep}{
	todo
}

\exe{difficulty=2}{
	Soit $d\in\N$. 
	Montrer que la famille de fonctions polynomiales sur $\R$ $\{1 ; x ; \dots ; x^d\}$ est linéairement indépendante sur $\R$.
}{exe:1xxd-indep}{
	todo
}

\exe{difficulty=1}{
	Soit $S = \{ -1 ; 0 ; 1\}$.
	Considérons l'ensemble $S[x]$, les polynômes à coefficients dans $S$.
	Par définition, le polynôme $x^3 - x$ n'est pas nul.
	
	Montrer que la fonction polynomiale $f(x) = x^3 -x$, elle, est la fonction nulle dans $S[x]$.
}{exe:fonction-nulle-S}{
	Il s'agit de montrer que $f(x) = 0$ pour tout $x\in S$.
	Ceci se vérifie par le calcul : $f(0) = f(1) = f(-1) = 0$.
}

\exe{}{
	Montrer que la famille de polynômes $\{ x+1 ; x+2 \} \subseteq\R[x]$ est linéairement indépendante.
}{exe:lin-indep1}{
	todo
}

\exe{}{
	Montrer que la famille de polynômes $\{ x^2+4x - 2 ; x+2 ; 50 \} \subseteq\R[x]$ est linéairement indépendante.
}{exe:lin-indep2}{
	todo
}

\exe{}{
	Est-ce que la famille de polynômes $\{ 2x^2+4x - 2 ; x^2 + x + 1 ; x-2 \} \subseteq\R[x]$ est linéairement indépendante ?
	Démontrer que oui ou fournir un certificat de non indépendance linéaire : une combinaison linéaire nulle et non triviale.
}{exe:lin-indep3}{
	La famille est liée car
		\[ (2x^2 + 4x - 2) - 2(x^2 + x + 1) - 2(x-2) = 0. \]
}

\exe{difficulty=1}{
	Soit $d\in\N$ pas un carré parfait. 
	Montrer que la famille $\{ 1 ; \sqrt{d} \} \subseteq\R$ est linéairement indépendante sur $\Q$.
}{exe:lin-indep4}{
	todo
}


\prop{}{
	Considérons une famille $\{ f_1, f_2, \dots, f_n\} \subseteq \R[x]$ linéairement indépendante.
	Alors si $p\in\R[x]$ est combinaison linéaire des $f_i$, cette combinaison est unique.
}{prop:indep-unicite}

\pf{}{
	Prenons $p = \alpha_i f_i$ et $p = \beta_i f_i$ deux combinaisons linéaires de $p$.
	Alors $p-p = 0 = (\alpha_i - \beta_i) f_i$ est la combinaison nulle.
	Par indépendance linéaire, c'est la combinaison triviale et donc $\alpha_i - \beta_i = 0 \iff \alpha_i = \beta_i$ pour tout $i=1, 2, \dots, n$.
}

\section{Divisibilité et division euclidienne}

\dfn{divisiblité}{
	Soient deux éléments $f, g \in \R[x]$ deux polynômes à coefficients réels.
	Le polynôme $f$ \emphindex{divise} $g$ dans $\R[x]$ dès que
		\[ g = f\cdot q, \]
	où $q\in \R[x]$.
}{dfn:divisibilite}

\notation{
	La relation « $f$ divise $g$ » est notée $f | g$.
}

\nomen{
	Un diviseur non constant est appelé un \emphindex{diviseur propre}.
}

\nt{
	Par convention, nous posons $\deg(0) = \minfty$.
	Cette convention est utile pour préserver le premier point de la proposition \ref{prop:deg-div} suivante dans le cas $g=0$.
}

\prop{}{
	Soient $f, g\in\R[x]$. Alors
	\begin{enumerate}
		\item
		$\deg(fg) = \deg(f) + \deg(g)$
		\item
		Si $f|g$ et $g\neq0$, alors $\deg(f) \leq \deg(g)$
	\end{enumerate}
}{prop:deg-div}

\pf{}{
	Si $f = 0$ ou $g=0$, alors il n'y a rien à démontrer.
	
	Sinon, notons $f(x) = f_i x^i$ et $g(x) = g_i x^i$ de degrés $m$ et $n$ respectivement.
	Alors $f_m \neq 0$ et $g_n \neq 0$.
	Ainsi, $(fg)(x)= f_mg_nx^{m+n} + \dots$ et $\deg(fg) = m+n$ car $f_m g_n \neq 0$.
}


\exe{}{
	Quand est-ce que l'identité $\deg(f+g) = \max\{\deg(f) ; \deg(g)\}$ est vraie ?
}{exe:deg-sum}{
	Il faut et il suffit que la somme des coefficients de degré $\max\{\deg(f) ; \deg(g)\}$ soit non nulle.
	Si $g=-f$ par exemple, alors $\deg(f+g) = \minfty$.
}


\thm{division euclidienne}{
	Soient $f, g\in\R[x]$ deux polynômes à coefficients réels avec $g\neq0$.
	Il existe deux uniques polynômes $q, r\in\R[x]$ vérifiant
	\begin{enumerate}
		\item $f = bg + r$ ; et
		\item $\deg(r) < \deg(g)$
	\end{enumerate}
	C'est la \emphindex{division euclidienne} de $f$ par $g$.
}{thm:div-euclidienne}

\nt{
	Il suit naturellement que $f$ divise $g$ dans $\R[x]$ si et seulement si le reste $r$ de la division euclidienne de $f$ par $g$ est nul.
	
	Remarquons que si $g=0$, alors la division euclidienne de $f$ par $g$ soit n'existe pas (si $f\neq0$), soit n'est pas unique (si $f=0$).
}

\pf{}{\leavevmode

	\underline{Existence.}

	Si $\deg(f) < \deg(g)$, alors il suffit de prendre $q=0$ et $r=g$.
	Sinon, $\deg(f) \geq \deg(g)$ et nous construisons $q$ en annulant le degré le plus haut de $f$ itérativement.
	
	\begin{enumerate}
		\item
		Initialiser $q = 0$.
		\item
		Tant que $\deg(f-gq) \geq \deg(g)$, 
		\item
		Ajouter un multiple de $x^{\deg(f-gq) - \deg(g)}$ à $q$ de telle sorte que le degré de $f-gq$ diminue d'au moins 1.
		\item
		Fin Tant que
		\item
		Renvoyer $q$ et $r=f-gq$.
	\end{enumerate}

	\underline{Unicité.}
	
	Si $f = gq + r = gq' + r'$ avec $\deg(r), \deg(r') < \deg(g)$, alors
		\[ g(q-q') = r' - r, \]
	et donc $q-q'=0$ car $\deg(r'-r) < \deg(g)$ et que $g\neq0$.
	Il suit que $q=q'$ et donc que $r=r'$.
}

\exe{}{
	Considérons l'équation quadratique $3x^2-7x+4= 0$ pour $x\in\R$.
	\begin{enumerate}
		\item
		Montrer que $1$ est racine de l'équation.
	\end{enumerate}
	Le but des questions suivantes est d'effectuer la division euclidienne de $3x^2-7x+4$ par $x-1$ afin de trouver la deuxième racine de l'équation.
	\begin{enumerate}[resume]
		\item
		Quelle fonction affine $q_1(x)$ prendre pour que $3x^2 -7x+4 - (x-1)q_1(x)$ n'ait plus de terme en $x^2$ ?
		\item
		Quelle fonction affine $q_2(x)$ prendre pour que $3x^2 -7x+4 - (x-1)q_1(x) - (x-1)q_2(x)$ n'ait plus de terme en $x$ ?
		\item
		Montrer que $3x^2 - 7x+4 = (x-1)\bigl[q_1(x) + q_2(x)\bigr]$ et en déduire l'autre racine de l'équation.
	\end{enumerate}
}{exe:euclide-poly}{
	todo
}


\section{Racines et multiplicités}

\dfn{racine, zéro}{
	Soit $f\in\R[x]$ un polynôme à coefficients réels.
	Un nombre réel $r \in \R$ est une \emphindex{racine} du polynôme $f$ ou \emphindex{zéro} de la fonction polynomiale $f$ dès que $f(r) = 0$.
}{dfn:racine}

\exe{}{
	On souhaite connaître les solutions de l'équation du deuxième degré d'inconnue $x\in\R$ :
		\begin{align}
			22x^2 - 125x + 22  = 0. \label{eq:1}
		\end{align}
	\begin{enumerate}
		\item Montrer que $22x^2 - 125x + 22 = (2x-11)(11x-2)$.
		\item En déduire l'ensemble des solutions de l'équation \eqref{eq:1}.
	\end{enumerate}
}{exe:7}{
	\begin{enumerate}
		\item On part toujours de la forme factorisée qu'on développe par double distributivité :
			\[ (2x - 11)(11x - 2) = 22x^2 - 4x - 121x + 22 = 22x^2 - 125x + 22. \]
		\item On utilise que le produit est nul si et seulement si un des deux facteurs est nul.
		L'ensemble des solutions de \eqref{eq:1} est donc $\left\{ \dfrac{11}2 ; \dfrac2{11} \right\}$.
	\end{enumerate}
}

\exe{}{
	Soit $f(x) = x^2 - x - 1$ une fonction quadratique sur $\R$.
	
	\begin{enumerate}
		\item
		Montrer que  $f(x) = \left( x - \frac{1 - \sqrt{5}}2 \right)\cdot \left( x - \frac{1 + \sqrt{5}}2 \right)$ pour tout $x\in\R$.
		\item 
		En déduire les zéros de $f$. Lequel est positif et lequel est négatif ? Raisonner sans calculatrice.
	\end{enumerate}
	Le zéro positif $\phi$ (lu \og phi \fg) s'appelle le \emph{nombre d'or}.
	\begin{enumerate}[resume]
		\item
		Vérifier que les zéros trouvés annulent bien $f$ en calculant leur image par $f$.
	\end{enumerate}
}{exe:phi}{
	todo
}

\exe{}{
	Pour chaque propriété, donner une fonction polynomiale $f$ sur $\R$ non identiquement nulle la vérifiant.
	\begin{multicols}{2}
	\begin{enumerate}
		\item $f$ s'annule en $1$.
		\item $f$ s'annule en $-10$.
		\item $f$ s'annule en $0$.
		\item $f$ s'annule en $-10$ et en $1$.
		\item Les zéros de $f$ sont $2, -3, \frac27$, et $0$.
	\end{enumerate}
	\end{multicols}
}{exe:6}{
	Le polynôme identiquement nul $f(x) = 0$ n'est bien sûr pas autorisé sinon l'exercice n'a pas d'intérêt !
	\begin{enumerate}
		\item $f(x) = x-1$ fonctionne ainsi que $f(x) = 3(x-1)$, ou $f(x) = -(x-1) = 1-x$.
		\item $f(x) = x+10$ fonctionne, ainsi que $f(x) = (x+10)(x-1)$, et tous ses multiples.
		\item $f(x) = x$ ou $f(x) = 500x$.
		\item $f(x) = (x+10)(x-1) = x^2 +9x -10$.
		\item $f(x) = (x-2)(x+3)(x-\frac27)x$.
	\end{enumerate}

}

\prop{}{
	Soit $f\in\R[x]$ un polynôme non nul à coefficient réels et $r\in\R$ une racine.
	Alors le monôme $x-r$ divise $f$ dans $\R[x]$ : il existe un $q\in\R[x]$ non nul vérifiant
		\[ f(x) = (x-r) q(x). \]
}{prop:x-r-div}

\pf{}{
	Soit $f(x) = (x-r)q(x) + c(x)$, $\deg(c) < \deg(x-r)$ la division euclidienne de $x-r$ par $f$.
	\begin{enumerate}
		\item
		La condition des degrés implique que $c$ est une constante.
		\item
		Le caractère non nul de $f$ implique que $q$ est non nul.
		\item
		L'évaluation en $x=r$ aboutit à $f(r) = q(r) (r-r) + c \iff 0 = c$.
	\end{enumerate}
	Par conséquent, $f(x) = (x-r) q(x)$.
}


\exe{}{
	Trouver une racine triviale puis factoriser le polynôme $f(x) = 7x^2 - 10x + 3$ pour trouver sa deuxième racine.
}{exe:euclide-poly2}{
	todo
}


\exe{}{
	Considérons le polynôme $f(x) = 21x^3 - 7x^2 - 9x + 3$.
	Montrer, par le calcul, que $\frac13$ est une racine de $f$.
	En déduire les autres racines de $f$.
	
	\emph{On pourra calculer la division euclidienne de $f$ par $3x-1$ plutôt que par $x-\frac13$.}
}{exe:euclide-poly3}{
	$f(x) = (3x-1)(7x^2-3)$ donc les racines sont $\frac13$ et $\pm\sqrt{\frac37}$.
}

%
%\thm{fondamental de l'algèbre}{
%	Soit $f\in\R[x]$ un polynôme non constant à coefficients réels.
%	Alors $f$ admet une racine dans $\C$.
%}{thm:fond-algebre}
%
%\exe{}{
%	Supposons que $z\in\C$ soit racine du polynôme $f\in\R[x]$ à coefficients réels.
%	Montrer que son conjugué complexe $\overline{z}$ est également racine.
%}{exe:Galois2}{
%	todo
%}
%
%\lem{}{
%	Soit $f\in\R[x]$ un polynôme non constant à coefficients réels et $z\in\C$ une de ses racines complexes.
%	Alors le polynôme $(x-z)(x-\overline{z})$ est un polynôme à coefficients réels qui divise $f$.
%}{lem:conjugate-roots}
%
%
%\thm{}{
%	Soit $f\in\R[x]$ un polynôme à coefficients réels et $r\in\C$ une racine.
%	Alors $(x-r)$ divise $f$ dans $\C[x]$.
%	Si $r\in\R$ est réelle, alors $(x-r)$ divise $f$ dans $\R[x]$.
%}{thm:facteur-lin}

\dfn{multiplicité}{
	Soit $f\in\R[x]$ un polynôme à coefficients réels et $r\in\R$ une racine.
	La \emphindex{multiplicité} $m$ de la racine $r$ est l'entier tel que $(x-r)^{m}$ divise $f$ dans $\R[x]$ mais pas $(x-r)^{m+1}$.
}{dfn:multiplicite}

\nomen{
	Une racine est dite \emphindex{racine simple} si sa multiplité est égale à $1$.
}

\cor{}{
	Soit $f\in\R[x]$ un polynôme à coefficients réels.
	Alors $f$ se \emphindex{factorise} comme
		\[ f(x) = \left(\prod_{\text{$r$ racine de $f$}} (x-r)^{m(r)}\right) \cdot q(x), \]
	où $q(x)$ n'admet aucune racine réelle.
}{cor:scinde}

\exe{}{
	Montrer qu'une fonction polynomiale de $\R[x]$ de degré 3 admet au moins une racine réelle.
}{exe:deg3-root}{
	todo
}

\section{Dérivation formelle}

\nomen{
	Un \emphindex{opérateur} est un autre nom pour une fonction. 
	Typiquement, un opérateur ne prend pas un nombre mais un objet plus complexe, comme un vecteur ou une fonction.
	Nous considérons dans ce qui suit l'opérateur dérivé qui prendre un polynôme et rend son polynôme dérivé.
	
	Lorsqu'une famille $\{ p_1  ; p_2 ; \dots ; p_k \}$ est linéairement indépendante, il est possible de définir un opérateur $T$ uniquement sur ses éléments pour le définir sur toutes les combinaisons linéaires par
		\[ T\left(\alpha_1p_1 + \alpha_2 p_2 + \dots + \alpha_k p_k\right) = \alpha_1T\left(p_1\right)  + \alpha_2T(p_2)  + \dots + \alpha_k T\left(p_k\right). \]
	La proposition \ref{prop:indep-unicite} d'unicité des coefficients $\alpha_i$ est clef ici.
	Sans celle-ci, deux combinaisons distinctes pourraient aboutir à deux images différentes d'un même antécédent.
}

\dfn{dérivation formelle}{
	L'\emphindex{opérateur dérivé} est noté $(\cdot)'$ et lu « prime » vérifie
		\begin{enumerate}
			\item 
			$(1)' = 0$ ;
			\item 
			$(x^n)' = nx^{n-1}$ pour $n\in\N^*$ ; et
			\item 
			est \emphindex{étendu linéairement} par $(\alpha f + g)' = \alpha f' + g'$ pour $\alpha\in\R$ et $f, g \in \R[x]$.
		\end{enumerate}
	Autrement dit, le \emphindex{polynôme dérivé} est défini par
		\[ (a_0 + a_1x + a_2x^2 + \dots + a_d x^d)' = a_1 + 2a_2x + 3a_3x^2 + \dots + da_d x^{d-1}. \]
}{dfn:derivation}

\notation{
	Nous notons la double dérivée $(f')'$ par $f''$ et la $k$-ième dérivée par $f^{(k)}$.
	Ainsi, $f^{(0)} = f$ et $f^{(k+1)} = \left(f^{(k)}\right)'$ pour $k\in\N$.
}

\notation{
	La \emphindex{factorielle} d'un entier naturel $n\in\N$, notée $n!$ et lue « $n$ factorielle » ou « factorielle $n$ » est définie par
	\begin{align*}
		0! = 1 && \et && n! = n \times (n-1) \times (n-2) \times \cdots \times 2 \times 1.
	\end{align*}
}

\exe{}{
	Soient $k, n\in\N$.
	Montrer que la $k$-ième dérivée de $x^n$ est égale à $0$ si $n<k$ et $k! x^{n-k}$ sinon.
}{exe:deriv-xn}{
	todo
}

\exe{difficulty=1}{
	Montrer la \emphindex{règle du produit} : pour deux polynômes $f, g\in\R[x]$,
		\[ (fg)' = f'g + g'f. \]
}{exe:deriv-prod}{
	À gauche, le coefficient en $x^n$ de $fg$ est égal à
		\[ \sum_{i+j=n} f_i g_j. \]
	Le coefficient en $x^{n-1}$ de $fg$ est donc égal à
		\[ n\sum_{i+j=n} f_i g_j. \]
	À droite, le coefficient en $x^{n-1}$ de $f'g$ est égal à
		\[ \sum_{i+j=n-1} f'_ig_j = \sum_{i+j=n-1} (i+1)f_{i+1}g_j = \sum_{i+j=n} if_ig_j. \]
	Le coefficient en $x^{n-1}$ de $g'f$ est égal à
		\[ \sum_{i+j=n-1} f_ig'_j = \sum_{i+j=n-1} f_i (j+1)g_{j+1} = \sum_{i+j=n} f_ijg_j. \]
	Le coefficient en $x^{n-1}$ de $f'g + g'f$ est donc égal à
		\[ \sum_{i+j=n} f_ijg_j +  \sum_{i+j=n} if_ig_j = \sum_{i+j=n} (i+j)f_ig_j = n\sum_{i+j=n} f_i g_j. \]
	Les coefficients de $(fg)'$ et $f'g+g'f$ sont donc égaux et les polynômes sont égaux.
}

\thm{de Taylor polynomial}{
	Soit $f\in\R[x]$ un polynôme à coefficients réels de degré $d$.
	Alors
		\[ f(x) = \sum_{k=0}^{d} \frac{f^{(k)}(0)}{k!} x^k.\]
}{thm:taylor-poly}

\pf{}{
	Posons $f(x) = a_0 + a_1 x + \dots + a_d x^d$.
	Il suffit de montrer que $a_k = \frac{f^{(k)}(0)}{k!}$ pour chaque $k=0, 1, 2, \dots, d$.
	Or, d'après l'exercice \ref{exe:deriv-xn} et par linéarité de la dérivation, nous avons
		\[ f^{(k)} = k! \left( a_k + a_{k+1} x + a_{k+2} x^2 + \dots + a_d x^{d-k} \right). \]
	Par conséquent, $f^{(k)}(0) = k! \cdot a_k$, et diviser par $k!$ conclut.
}

\nt{
	Une connaissance locale d'une fonction polynomiale et de ses dérivées permet donc de la connaître partout.
	Le \emphindex{théorème de Taylor} se généralise : sous certaines conditions --- moins fortes qu'être polynomiale --- nous avons
		\[ f(x) = \sum_{k=0}^{d} \frac{f^{(k)}(0)}{k!} x^k + o(x^d), \]
	où $o(x^d)$ désigne une fonction dominée par $x^d$ autour de 0, comme par exemple $x^{d+1}$ dont le ratio $\frac{x^{d+1}}{x^d}$ tend vers 0 lors que $x$ tend vers 0.
}

\exe{difficulty=1}{
	Soit $f\in C^1(\R)$ une fonction continûment dérivable.
	Montrer que 
		\[ f(x) = f(0) + f'(0)x + g(x), \]
	où $g(x) = o(x)$ dans le sens suivant : 
		\[ \lim_{x \to 0} \frac{g(x)}x = 0. \]
}{exe:taylor-C1}{
	Posons $g(x) = f(x) - f(0) - f'(0)x$.
	Pour $x\neq0$, nous avons
		\[ \frac{g(x)}x = \frac{f(x) - f(0)}{x-0} - f'(0). \]
	Lorsque $x$ tend vers $0$,  $\frac{f(x) - f(0)}{x-0}$ tend vers $f'(0)$ par définition de $f'(0)$.
	Il suit donc que $\frac{g(x)}x$ tend vers 0 comme requis.
	
	Remarque : si $f\in C^2(\R)$, il s'avère que $g(x) = \O(x^2)$, c'est-à-dire que $\frac{g(x)}{x^2}$ reste borné lorsque $x$ tend vers $0$ (ce qui est plus fort que $g=o(x)$).
	Nous aurons en fait 
		\[ f(x) = f(0) + f'(0)x + \frac{f''(0)}2 x^2 + O(x^2). \]
}

\exe{difficulty=1}{
	Montrer le théorème de Taylor bis : 
	soit $f\in\R[x]$ un polynôme à coefficients réels de degré $d$ et $a\in\R$ un réel quelconque.
	Alors
		\[ f(x) = \sum_{k=0}^{d} \frac{f^{(k)}(a)}{k!} (x-a)^k.\]
}{exe:taylor-poly-2}{
	todo
}

\exe{}{
	Considérons le polynôme
		\[ f(x) = (x-3)^2 - (x-3)^4 + 3(x-3)^5. \]
	Calculer $f(3), f'(3), f''(3), f'''(3), f^{(\text{iv})}(3)$, et $f^{(\text{v})}(3)$.
}{exe:f-shift-derivatives}{

}

\prop{}{
	Soit $f\in\R[x]$ un polynôme à coefficients réels de degré $d$ et $a\in\R$ un réel quelconque.
	Montrer que $a$ est une racine de multiplicité $m$ si et seulement si $f(a) = f'(a) = \cdots = f^{(m-1)}(a) = 0$ et $f^{(m)}(a) \neq 0$.
}{prop:mult-from-diff}

\exe{difficulty=1}{
	Démontrer la proposition \ref{prop:mult-from-diff}.
}{exe:mult-from-diff}{
	todo
}

\thm{}{
	Soit $f\in\R[x]$ un polynôme à coefficients réels. 
	Supposons de sucroît que tous les racines de $f$ soient réelles\footnote{hypothèse dont on pourra se passer plus tard grâce aux nombres complexes.}.
	
	Alors les seuls diviseurs d'à la fois $f$ et $f'$ sont les constantes si et seulement si toutes les racines de $f$ sont simples.
}{thm:derivative-test}

\exe{difficulty=1}{
	Démontrer le théorème \ref{thm:derivative-test}.
}{exe:derivative-test}{
	\underline{Direction $\implies$.}
	
	Soit $r\in\R$ une racine de $f$ de multiplicité $m\in\N^*$.
	Comme $x-r$ divise $f$ mais pas $f'$, nous avons $f(r) = 0$ et $f'(r) \neq 0$.
	La proposition \ref{prop:mult-from-diff} implique que $m=1$.
	
	\underline{Direction $\impliedby$.}
	
	Montrons la contraposée.
	Soit $d\in\R[x]$ un diviseur non constant à la fois de $f$ et de $f'$.
	Les racines de $d$ étant des racines de $f$, elles sont toutes réelles.
	Si $d$ n'étant pas constant, il existe un réel $r\in\R$ tel que $x-r$ divise $d$ et donc $f$ et $f'$.
	Ainsi $f(r) = f'(r) = 0$ et la racine $r$ n'est pas simple.
}



\exe{difficulty=1}{
	Considérons le \emphindex{symbole de Pochhammer} défini par
		\begin{align*}
			(x)_0 = 1, && \et && (x)_n = x(x-1)(x-2)\cdots(x-n+1) \text{ pour $n\in\N^*$, }
		\end{align*}
	ainsi que l'opérateur de différence $\Delta$ qui associe à chaque $f\in\R[x]$ le polynôme $f(x+1)-f(x)$.
	
	Montrer que, pour tout $n\in\N^*$, l'identité suivante est vraie dans $\R[x]$.
		\[ \Delta(x)_n = n(x)_{n-1} \]
}{exe:Pochhammer}{

}

\section*{Exercices supplémentaires}
\addcontentsline{toc}{section}{Exercices supplémentaires}

\exe{}{
	On souhaite connaître les solutions de l'équation du troisième degré d'inconnue $x\in\R$ :
		\begin{align}
			4x^3 - 6x^2 - 2x + 3  = 0. \label{eq:2}
		\end{align}
	\begin{enumerate}
		\item Montrer que $4x^3 - 6x^2 - 2x + 3 = (2x-3) \left(2x^2-1\right)$.
		\item En déduire l'ensemble des solutions de l'équation \eqref{eq:2}.
	\end{enumerate}
}{exe:8}{

	\begin{enumerate}
		\item On part toujours de la forme factorisée qu'on développe par double distributivité.
		On utilise aussi que $x^2 \cdot x = x \cdot x \cdot x = x^3$.
			\[ (2x - 3)(2x^2 - 1) = 4x^3 - 2x -6x^2 + 3. \]
		\item On utilise que le produit est nul si et seulement si un des deux facteurs est nul.
		Or l'équation $2x^2 - 1 = 0$ est équivalente à $x^2 = \dfrac12$.
		En prenant la racine et en utilisant que $\sqrt{x^2} = |x|$, on obtient
			\[ |x| = \sqrt{\dfrac12} = \dfrac1{\sqrt2}. \]
		Les deux valeurs $\pm \dfrac1{\sqrt2}$ (notation $\pm$ pour \og plus ou moins \fg) sont donc solutions.
		On en déduit que l'ensemble des solutions de \eqref{eq:2} est $\left\{ \dfrac32 ; \pm \dfrac1{\sqrt2} \right\} = \left\{ \dfrac32 ; \dfrac1{\sqrt2} ; - \dfrac1{\sqrt2} \right\}$
	\end{enumerate}

}


\exe{}{
	On souhaite connaître les solutions de l'équation du quatrième degré d'inconnue $x\in\R$ :
		\begin{align}
			16x^4 - 24x^2 + 9  = 0. \label{eq:3}
		\end{align}
	\begin{enumerate}
		\item Montrer que $16x^4 - 24x^2 + 9 = \left(4x^2-3\right)^2$.
		\item En déduire l'ensemble des solutions de l'équation \eqref{eq:3}.
	\end{enumerate}
}{exe:9}{

	\begin{enumerate}
		\item On part toujours de la forme factorisée qu'on développe à l'aide de l'identité remarquable $(a-b)^2 = a^2 + b^2 -2ab$.
		On utilise aussi que $(x^2)^2 = x^2 \cdot x^2 = x\cdot x\cdot x\cdot x = x^4$.
			\[ (4x^2 - 3)^2 = (4x^2)^2 + 3^2 - 2\cdot(4x^2)\cdot3 = 16x^4 + 9 -24x^2. \]
		\item On utilise que le produit est nul si et seulement si un des deux facteurs est nul.
		Ici, les deux facteurs sont les mêmes, car $(4x^2 - 3)^2 = (4x^2 - 3)(4x^2 - 3)$, donc on se remet à résoudre $(4x^2 - 3) = 0$.
		Ceci est équivalent à $x^2 = \dfrac34$, et donc l'ensemble des $x$ vérifiant \eqref{eq:3} est donné par $\left\{ \pm \sqrt{\dfrac34} \right\} = \left\{ \pm \dfrac12\sqrt{3} \right\}= \left\{ \dfrac12\sqrt{3} ;  -\dfrac12\sqrt{3} \right\}$.
	\end{enumerate}

}





\exe{}{
	Pour chaque proposition suivante, démontrer qu'elle est \underline{toujours vraie} ou montrer qu'elle \underline{peut être fausse} avec un contre-exemple (un dessin peut suffir).
	
	$f, g, h, F$ et $G$ sont des fonctions définies sur $\R$.
	\begin{enumerate}
		\item Si $f$ est affine et admet deux zéros distincts, alors $f(x) = 0$ pour tout $x\in\R$.
		\item Si $g(x) \geq 0$ pour tout $x\in\R$, alors $g$ n'admet aucun zéro.
		\item Si $h(x) > 0$ pour tout $x\in\R$, alors $h$ n'admet aucun zéro.
		\item Si $F$ n'admet aucun zéro, alors $F(x) > 0$ pour tout $x\in\R$.
		\item $G(x) = (x-3)^2 (x-5)^2$ est toujours positive est admet exactement deux zéros distincts.
	\end{enumerate}
}{exe:14}{

	\begin{enumerate}
		\item 
		C'est vrai. 
		Si $f$ est affine et admet deux racines distinctes, alors sa courbe représentative est une droite qui intersecte deux fois l'axe des abscisses en deux endroits différents.
		$\C_f$ et l'axe des abscisses sont donc confondues et on a bien $f(x) = 0$ pour tout $x\in\R$.
		
		Algébriquement, on a $(x_A ; 0), (x_B ; 0) \in \C_f$ pour $x_A \neq x_B$.
		La formule du coefficient directeur donne $a = 0$, et une appartenance donne $b=0$.
		
		\item 
		C'est faux : être positif ou nul n'empêche pas d'être nul.
		On pourra prend la fonction carré $f(x) = x^2$ comme contre-exemple : $f(x) \geq 0$ pour tout $x\in\R$ et pourtant $f$ admet $0$ pour racine car $f(0) = 0$.
		
		\item 
		C'est vrai : être strictement positif empêche d'être nul. 
		C'est bien la différence entre la positivité stricte ($>$) et large ($\geq$).
		
		\item 
		C'est faux : si $F$ ne s'annule jamais, elle n'est pas nécessairement toujours strictement positive car elle pourrait aussi être toujours strictement négative !
		Prenons par exemple $f(x) = -1 -x^2$ qui vérifie $f(x) \leq -1 < 0$ pour tout $x\in\R$ réel.
		
		\item 
		C'est vrai.
		Un carré est toujours positif, donc $G(x)$ est le produit de deux positifs et est positif.
		En outre, si $G(x) = 0$, alors $(x-3)^2 = 0$ ou $(x-5)^2 = 0$.
		Il suit que $x=3$ ou $x=5$, et ce sont les deux seules racines de $G$.
		
	\end{enumerate}

}

\exe{difficulty=1}{
	Considérons une famille de polynômes $\{ f_1, f_2, \dots, f_n\} \subseteq \R[x]$ telle que pour tout $i=1, 2, \dots n$, il existe un nombre $r_i\in\R$ qui est racine de $f_j$ pour tout $j\neq i$, mais n'est pas racine de $f_i$.
	Montrer que la famille est linéairement indépendante.
	
	Conclure que la famille $F$ suivante est linéairement indépendante sur $\R$.
		\[ F = \bigset{ (x-1)(x-2)(x-3) ; x(x-2)(x-3) ; x(x-1)(x-3) ; x(x-1)(x-2) } \]
}{exe:unique-roots}{
	todo
}

\exe{difficulty=2}{
	Supposons que $r+s\sqrt{d}$, avec $r, s\in\Q$ et $d\in\N$ non carré parfait, soit racine de l'équation $ax^2 + bx + c = 0$ où $a, b, c\in\Q$.
	Montrer que son conjugué $r-s\sqrt{d}$ est également racine.
}{exe:Galois2}{
	Voici une première stratégie de preuve possible.
	\begin{enumerate}
		\item
		Montrer que $1$ et $\sqrt{d}$ sont \emph{linéairement indépendants} sur $\Q$.
		C'est-à-dire, montrer que si $r + s\sqrt{d} = 0$ avec $r, s \in\Q$, alors $r=s=0$.
		Une équation sur les réels devient alors deux équations sur les rationnels !
		\item
		En déduire que $a(r+s\sqrt{d})^2 + b(r+s\sqrt{d}) + c=0$ est équivalente aux deux équations $ar^2 + ads^2+br+c=0$ et  $2ars+bs=0$.
		\item
		Conclure en calculant $a(r-s\sqrt{d})^2 + b(r-s\sqrt{d}) + c$.
	\end{enumerate}
	À l'étape 2, nous utilisons aussi que la somme et le produit de deux rationnels est un rationnel, et c'est à ce moment que nous utilisons que tous les coefficients ($r, s, a, b, c$) sont rationnels.
	Cette propriété de $\Q$ se démontre assez simplement avec les formules $\frac{a}b+\frac{c}d$ et $\frac{a}b\times\frac{c}d$.
	
	Voici un deuxième stratégie de preuve possible.
	\begin{enumerate}
		\item
		Montrer que $1$ et $\sqrt{d}$ sont \emph{linéairement indépendants} sur $\Q$.
		C'est-à-dire, montrer que si $r + s\sqrt{d} = 0$ avec $r, s \in\Q$, alors $r=s=0$.
		\item
		Montrer que la fonction « conjugué » $f\bigl(r+s\sqrt{d}\bigr) = r - s\sqrt{d}$ où $r, s \in \Q$ vérifie que
			\begin{enumerate}[label=(\roman*)]
				\item $f(x+y) = f(x)+f(y)$
				\item $f(xy) = f(x)f(y)$
			\end{enumerate}
		\item
		Conclure en appliquant $f$ à l'équation $ax^2 +bx + c = 0$ pour obtenir $af(x)^2 + bf(x) + c = 0$.
	\end{enumerate}
	
}

\shipoutExercise

\section*{Solutions}
\addcontentsline{toc}{section}{Solutions}

\shipoutAnswer



